% Part: sets-functions-relations
% Chapter: size-of-sets-complete

\documentclass[../../../include/open-logic-chapter]{subfiles}

\begin{document}

\olchapter{sfr}{siz}{Hoe groot zijn verzamelingen?}

\begin{editorial}
Dit hoofdstuk gaat over opsommingen en over verzamelingen die wel of
niet aftelbaar zijn. Bij een aantal paragrafen zijn er twee versies. De
eenvoudigere versie ziet een opsomming als een lijst, of als een
surjectie vanuit $\PosInt$---een functie die alle gewenste elementen
bereikt. De abstractere versie ziet een bijectie met $\Nat$ als een
opsomming. Zo'n bijectie is een functie die niets dubbel telt en elk
element bereikt.
\end{editorial}

\olimport{introduction}

\olimport{enumerability}

\olimport{zig-zag}

\olimport{pairing}

\olimport{pairing-alt}

\olimport{non-enumerability}

\olimport{reduction}

\olimport{equinumerous-sets}

\olimport{comparing-size}

\olimport{schroder-bernstein}

\begin{editorial}
De drie volgende paragrafen zijn \olref[sfr][siz][enm-alt]{sec},
\olref[sfr][siz][nen-alt]{sec} en \olref[sfr][siz][red-alt]{sec}. Het
zijn andere versies van \olref[sfr][siz][enm]{sec},
\olref[sfr][siz][nen]{sec} en \olref[sfr][siz][red]{sec}. Tim Button
schreef ze voor zijn tekst Open Set Theory. Deze versies zijn iets
moeilijker. Ze gebruiken een andere betekenis van opsombaarheid die
beter bij verzamelingenleer past. Hier is iets opsombaar als er een
bijectie is met $\Nat$ of met een beginsegment, dus een eerste stuk.
Dat is iets anders dan zeggen dat je de elementen op een lijst kunt
zetten, of dat ze precies de uitkomsten zijn van een surjectieve functie
met $\PosInt$ als domein.
\end{editorial}

\olimport{enumerability-alt}
\olimport{non-enumerability-alt}
\olimport{reduction-alt}

\OLEndChapterHook
\end{document}
